\chapter{Testing}
TBD
\begin{figure}[htp]
    \centering
    \scalebox{0.8}{
        \begin{tikzpicture}[node distance=2 cm]
            % \draw[help lines] (0, -3) grid (10,3);
            \airsupply{origin}
            \draw[thick, densely dotted] (1, -1.5) -| +(4, 3.5) -| ++(0, 0);
            \node (airsupplytext) at (3, -2) {Air Supply};

            \aircgg{airsupply}
            \draw[thick, densely dotted] (aircgg) ++(0.75, -1.5) -| +(-5.5, 3.5) -| ++(0, 0);
            \node (airsupplytext) at (8, -2) {CGG Air Equivalent};

            \node[pressure regulator, right of=aircgg] (pregulator) {};
            \draw[thick] (aircgg) -- (pregulator);

            \airfc{pregulator}
            \draw[thick, densely dotted] (airfc) ++(0.75, -1) -| +(-4.5, 8) -| ++(0, 0);
            \node (airsupplytext) at (15.5, -4) {Fuel Cell Air Equivalent};

        \end{tikzpicture}
    }
    \caption{Air Test Setup}%
    \label{fig:air test}%
\end{figure}



Goals of various tests:
\begin{itemize}
    \item Compressed Air Equivalent Test
          \begin{itemize}
              \item Explore unknown risks
              \item Verify expected system behaviour
              \item Verify simulation results
              \item Compare control strategies
          \end{itemize}

    \item FCU Pressure test
          \begin{itemize}
              \item Verify consistency of FCU output pressure
          \end{itemize}

    \item FCU mass flow test
          \begin{itemize}
              \item Verify consistency of FCU output mass flow
              \item Verify leakage of FCU
          \end{itemize}

    \item Gas supply test
          \begin{itemize}
              \item Verify consistency of gas pressure during purge events
              \item Verify consistency of gas mass flow during purge events
              \item Verify total leakage of gas supply system
          \end{itemize}

    \item Fuel Cell test
          \begin{itemize}
              \item Verify the fuel cell is working normally
              \item Verify the fuel cell output voltage
              \item Verify the fuel cell output power
              \item Verify the fuel cell fuel consumption
          \end{itemize}

    \item Propulsion test
          \begin{itemize}
              \item Verify the maximum thrust
              \item Verify the maximum power draw
              \item Verify the relation between input power and thrust output
          \end{itemize}

    \item Hybrid Battery test
          \begin{itemize}
              \item Verify the battery cannot back-feed power into the fuel cell
              \item Verify output power of hybrid system
              \item Verify output voltage of hybrid system
          \end{itemize}

    \item Integration test
          \begin{itemize}
              \item Verify all individually tested components behave as expected when integrated into one system
              \item Verify relation between fuel consumption and thrust output
          \end{itemize}
\end{itemize}

\clearpage
\section{Compressed Air Equivalent System}
TBD

Goals of compressed air equivalent tests:
\begin{itemize}
    \item Explore unknown risks
    \item Verify expected system behaviour
    \item Verify simulation results
    \item Compare control strategies
\end{itemize}

\subsection{Fuel Supply Regulator}
TBD

\begin{figure}[htp]
    \begin{subfigure}{0.5\textwidth}%
        \centering%
        \begin{tikzpicture}[node distance=2 cm]
            \node (in) [io] {Input: CGG Air Equivalent};
            \node (valve) [process, below of=in] {Stellar Space FCU};
            \node (out) [io, below of=valve] {Output: Fuel Cell Air Equivalent};

            \draw [arrow] (in) -- node[anchor=east] {0.35--3 Bar} (valve);
            \draw [arrow] (valve) -- node[anchor=east] {0.35--0.5 Bar} (out);
        \end{tikzpicture}%
        \caption{FCU}%
        \label{fc:fcu}%
    \end{subfigure}%
    \hfill%
    \begin{subfigure}{0.5\textwidth}%
        \centering%
        \begin{tikzpicture}[node distance=2 cm]
            \node (in) [io] {Input: CGG Air Equivalent};
            \node (valve) [process, below of=in] {Pressure Regulator};
            \node (out) [io, below of=valve] {Output: Fuel Cell Air Equivalent};

            \draw [arrow] (in) -- node[anchor=east] {0.35--3 Bar} (valve);
            \draw [arrow] (valve) -- node[anchor=east] {0.35--0.5 Bar} (out);
        \end{tikzpicture}%
        \caption{Pressure Regulator}%
        \label{fc:regulator}%
    \end{subfigure}%

    \caption{Compressed Air Equivalent}%
    \label{fc:air equivalent}%
\end{figure}%

\clearpage
\subsection{Cool Gas Generator}
TBD
\begin{figure}[htp]
    \begin{subfigure}{0.5\textwidth}%
        \centering%
        \begin{tikzpicture}[node distance=2 cm]
            \node (in1) [io] {Input: Shop Air};
            \node (pro1) [process, below of=in1] {Pressure Regulator};
            \node (pro2) [process, below of=pro1] {Ball Valve `Orifice'};
            \node (pro3) [process, below of=pro2] {Buffer Volume};
            \node (out1) [io, below of=pro3] {Output: Simulated CGG Production};

            \draw [arrow] (in1) -- node[anchor=east] {6--8 Bar} (pro1);
            \draw [arrow] (pro1) -- node[anchor=east] {6.0 Bar} (pro2);
            \draw [arrow] (pro2) -- node[anchor=east] {0.2 g/min} (pro3);
            \draw [arrow] (pro3) -- node[anchor=east] {3 Bar Max} (out1);
        \end{tikzpicture}%
        \caption{Ball Valve}%
        \label{fc:ball cgg}%
    \end{subfigure}%
    \hfill%
    \begin{subfigure}{0.5\textwidth}%
        \centering%
        \begin{tikzpicture}[node distance=2 cm]
            \node (in1) [io] {Input: Shop Air};
            \node (pro1) [process, below of=in1] {Pressure Regulator};
            \node (pro2) [process, below of=pro1] {Bronckhorst Mass Flow Controller};
            \node (pro3) [process, below of=pro2] {Buffer Volume};
            \node (out1) [io, below of=pro3] {Output: Simulated CGG Production};

            \draw [arrow] (in1) -- node[anchor=east] {6--8 Bar} (pro1);
            \draw [arrow] (pro1) -- node[anchor=east] {6.0 Bar} (pro2);
            \draw [arrow] (pro2) -- node[anchor=east] {0.2 g/min} (pro3);
            \draw [arrow] (pro3) -- node[anchor=east] {6 Bar Max} (out1);
        \end{tikzpicture}%
        \caption{Bronckhorst}%
        \label{fc:bronckhorst cgg}%
    \end{subfigure}%

    \caption{CGG Compressed Air Equivalent}%
    \label{fc:cgg equivalent}%
\end{figure}%

\clearpage
\subsection{Fuel Cell}
TBD
\begin{figure}[htp]
    % \centering%
    \begin{subfigure}{\textwidth}%
        \centering%
        \begin{tikzpicture}[node distance=2 cm]
            \node (in) [io] {Input: Regulated air from CGG equivalent};
            \node (valve) [decision, below of=in, yshift=-5mm] {5/2 Valve};
            \node (purge) [process, below of=valve, xshift = -3cm] {Purging: Ball Valve `Orifice'};
            \node (normal) [process, below of=valve, xshift = 3cm] {Normal Operation: Ball Valve `Orifice'};
            \node (outp) [io, below of=purge] {Output: Purged to atmosphere};
            \node (outn) [io, below of=normal] {Output: `Used by fuel cell' (released to atmosphere)};

            \draw [arrow] (in) -- node[anchor=east] {0.35--0.5 Bar} (valve);
            \draw [arrow] (valve) -| (purge);
            \draw [arrow] (valve) -| (normal);
            \draw [arrow] (purge) -- node[anchor=east] {0.6 g/min} (outp);
            \draw [arrow] (normal) -- node[anchor=east] {0.2 g/min} (outn);
        \end{tikzpicture}%
        \caption{Ball Valve, Single Switch}%
        \label{fc:ball fc}%
    \end{subfigure}%

    \vspace{1cm}
    \begin{subfigure}{\textwidth}%
        \centering%
        \begin{tikzpicture}[node distance=2 cm]
            \node (in) [io] {Input: Regulated air from CGG equivalent};
            \node (pvalve) [decision, below of=in, yshift=-5mm, xshift=-5cm] {5/2 Valve};
            \node (nvalve) [decision, below of=in, yshift=-5mm, xshift=5cm] {5/2 Valve};
            \node (porifice) [process, below of=pvalve] {Purging: Ball Valve `Orifice'};
            \node (norifice) [process, below of=nvalve] {Normal Operation: Ball Valve `Orifice'};
            \node (pout) [io, below of=porifice] {Output: Purged to atmosphere};
            \node (nout) [io, below of=norifice] {Output: `Used by fuel cell' (released to atmosphere)};
            \node (block) [io, below of= in, yshift=-4.5cm] {Output:\\Blocked};


            \draw [arrow] (in) -| node[anchor=east] {0.35--0.5 Bar} (pvalve);
            \draw [arrow] (in) -| node[anchor=west] {0.35--0.5 Bar} (nvalve);
            \draw [arrow] (pvalve) -- (porifice);
            \draw [arrow] (nvalve) -- (norifice);
            \draw [arrow] (pvalve) -| (block);
            \draw [arrow] (nvalve) -|    (block);
            \draw [arrow] (porifice) -- node[anchor=east] {0.6 g/min} (pout);
            \draw [arrow] (norifice) -- node[anchor=east] {0.2 g/min} (nout);
        \end{tikzpicture}%
        \caption{Ball Valve, Dual Switch}%
        \label{fc:ball fc}%
    \end{subfigure}%

    \vspace{1cm}
    \begin{subfigure}{\textwidth}%
        \centering%
        \begin{tikzpicture}[node distance=2 cm]
            \node (in) [io] {Input: Regulated air from CGG equivalent};
            \node (valve) [process, below of=in] {Bronckhorst Mass Flow Controller};
            \node (out) [io, below of=valve] {Output: H2 used or purged by Fuel Cell (released to atmosphere)};

            \draw [arrow] (in) -- node[anchor=east] {0.35--0.5 Bar} (valve);
            \draw [arrow] (valve) -- node[anchor=east] {Mass flow by Micro Controller} (out);
        \end{tikzpicture}%
        \caption{Bronckhorst}%
        \label{fc:bronckhorst fc}%
    \end{subfigure}%

    \caption{Fuel Cell Compressed Air Equivalent}%
    \label{fc:fc equivalent}%
\end{figure}%

\begin{tabularx}{\textwidth}{lXX}
    \toprule
    \textbf{Test setup}      & \textbf{Capabilities}              & \textbf{Downsides}                                         \\
    \midrule
    Ball Valve Single Switch & Test uncontrolled system behaviour & Mass flow is only accurate when input pressure is constant \\
    Ball Valve Dual Switch   & Test On--Off controlled behaviour  & Mass flow is only accurate when input pressure is constant \\
    Bronchorst Valve         & Test PID controlled behaviour      & Expensive                                                  \\
    \bottomrule
\end{tabularx}

\clearpage
\section{Hydrogen Test}
TBD
\begin{figure}[htp]
    \centering
    \begin{tikzpicture}[node distance=2 cm]
        \node (in1) [io] {Input: Hydrogen Gas Storage};
        \node (pro1) [process, below of=in1] {Pressure Reducer};
        \node (pro2) [sensor, below of=pro1] {Mass Flow Sensor};
        \node (pro3) [process, below of=pro2] {Flow Control Unit};
        % \node (pro4) [sensor, below of=pro3, yshift=-1cm] {Mass Flow Sensor};
        \node (pro5) [sensor, below of=pro3, yshift=0.5cm, xshift=6cm] {Pressure Sensor};
        \node (pro6) [process, below of=pro3, yshift=-1cm] {40 micron filter};
        \node (pro7) [process, below of=pro6] {Fuel Cell};
        \node (out1) [io, below of=pro7] {Output: 800W};

        \draw [arrow] (in1) -- node[anchor=east] {High pressure} (pro1);
        \draw [arrow] (pro1) -- node[anchor=east] {Low pressure} (pro2);
        \draw [arrow] (pro2) -- (pro3);
        \draw [arrow] (pro3) -- node[anchor=east] {0.9 Bar} (pro6);
        \draw [arrow] (pro3) |- (pro5);
        \draw [arrow] (pro5) |- (pro3);
        % \draw [arrow] (pro4) -- (pro6);
        \draw [arrow] (pro6) -- (pro7);
        \draw [arrow] (pro7) -- (out1);
    \end{tikzpicture}

    \caption{Hydrogen Test Setup}
    \label{fc:h2 test}
\end{figure}

\begin{figure}[htp]
    \centering
    \begin{tikzpicture}[node distance=2 cm]
        \node (in1) [io] {Input: 800W};
        \node (pro1) [sensor, below of=in1] {Amp Meter};
        \node (pro2) [process, below of=pro1] {(Ideal) Diode\\Current Regulator};
        \node (pro3) [process, below of=pro1, xshift=6cm, yshift=-1.5cm] {Battery};
        \node (pro4) [sensor, below of=pro3, yshift=-1cm] {Volt Meter};
        \node (pro5) [sensor, below of=pro2, yshift=-1cm] {Amp Meter};
        \node (pro6) [process, below of=pro5, yshift=-1cm] {ESC};
        \node (pro7) [process, below of=pro6] {Motor};
        \node (pro8) [process, below of=pro7] {Propeller};
        \node (pro9) [sensor, below of=pro6, xshift=6cm] {Force Sensor};
        \node (out1) [io, below of=pro8] {Output: Thrust};

        \draw [arrow] (in1) -- (pro1);
        \draw [arrow] (pro1) -- (pro2);
        \draw [arrow] (pro2) |- (pro3);
        \draw [arrow] (pro3) -| (pro5);
        \draw [arrow] (pro4) -| (pro6);
        \draw [arrow] (pro5) -- (pro6);
        \draw [arrow] (pro6) -- (pro7);
        \draw [arrow] (pro7) -- (pro8);
        \draw [arrow] (pro7) -- (pro9);
        \draw [arrow] (pro8) -- (out1);
    \end{tikzpicture}

    \caption{Electronics Test Setup}
    \label{fc:electric test}
\end{figure}


\newpage
